% ============================================================
% Paper 1: PhotoMedGemma Overview
% "PhotoMedGemma: A Photonic Chip Compiler for Deploying
%  Medical AI in Resource-Constrained Environments"
% ============================================================
\documentclass[12pt, a4paper]{article}

\usepackage{amsmath, amssymb, amsthm}
\usepackage{graphicx}
\usepackage{hyperref}
\usepackage{booktabs}
\usepackage{multirow}
\usepackage{array}
\usepackage{geometry}
\usepackage{xcolor}
\usepackage{listings}
\usepackage{caption}
\usepackage{subcaption}
\usepackage{cite}
\usepackage{microtype}

\geometry{margin=2.5cm}

\hypersetup{
  colorlinks=true,
  linkcolor=blue!70!black,
  citecolor=blue!70!black,
  urlcolor=blue!70!black,
}

\definecolor{codegreen}{rgb}{0,0.6,0}
\definecolor{codegray}{rgb}{0.5,0.5,0.5}
\lstset{
  basicstyle=\ttfamily\footnotesize,
  commentstyle=\color{codegreen},
  keywordstyle=\color{blue},
  numberstyle=\tiny\color{codegray},
  breaklines=true,
  frame=single,
  language=Python,
}

% ----------------------------------------------------------------
\title{\textbf{PhotoMedGemma:}\\[4pt]
       \large A Photonic Chip Compiler for Deploying\\
       Medical Artificial Intelligence in\\
       Resource-Constrained Environments}

\author{
  Linda H.\\[2pt]
  \small Independent Research\\[2pt]
  \small \href{https://github.com/yourusername/photomedgemma}{github.com/yourusername/photomedgemma}
}

\date{March 2026}
% ----------------------------------------------------------------

\begin{document}
\maketitle

\begin{abstract}
We present \textbf{PhotoMedGemma}, an open-source compiler and simulation
framework that translates the weights of MedGemma-4B-IT --- Google's
4-billion-parameter multimodal medical AI --- into phase settings for silicon
photonic Mach-Zehnder interferometer (MZI) meshes.
The central insight is that transformer inference is a \emph{static} operation:
once trained, model weights never change, so the computation can be permanently
encoded into the geometry of a photonic chip, converting matrix-vector
multiplication into passive light propagation.
Our compiler decomposes each weight matrix via truncated singular value
decomposition (SVD) and maps the resulting unitaries to Clements MZI meshes.
We validate the pipeline on real MedGemma activations captured from five
breast cancer PNG images processed on Kaggle.
Compilation achieves machine-precision fidelity ($\approx 7 \times 10^{-15}$)
at both rank-64 and full rank.
Rank-64 MZI meshes retain $23$--$92\%$ of the full weight energy
depending on the projection, requiring $r(r-1)/2 = 2{,}016$ MZIs per mesh
versus $2{,}096{,}128$ MZIs at full rank.
The photonic architecture projects $\approx 8\,\text{W}$ total system power
versus $\approx 300\,\text{W}$ for an A100 GPU --- a $37\times$ improvement.
All code, compiled phase maps, and experimental data are publicly available.
The ultimate goal is a sub-\$100 photonic chip that runs world-class medical AI
in hospitals throughout Africa and the global south with no cloud dependency,
no data privacy risk, and no energy crisis.
\end{abstract}

% ----------------------------------------------------------------
\section{Introduction}
\label{sec:intro}

Artificial intelligence has achieved clinical-grade performance in medical
imaging, diagnostic reasoning, and treatment planning.
Yet the hardware required to run such models is concentrated in wealthy
countries, behind cloud APIs, and inside data centers consuming hundreds of
megawatts.
In 2024, we worked with neurosurgeons at \textbf{Kenyatta National Hospital
(KNH)} in Nairobi, Kenya --- the largest referral hospital in East and Central
Africa.
The clinical need was unambiguous: AI-assisted triage for cerebral hemorrhage,
herniation, and neurological emergencies from CT scans.
MedGemma~\cite{medgemma2024} can perform all of these tasks.
It cannot, however, be deployed at KNH due to four barriers: data sovereignty
law prohibiting foreign data transfer; insufficient internet bandwidth for
real-time inference; prohibitive cloud API costs; and sub-second latency
requirements for emergency triage.

The answer is not a smaller model.
The answer is \emph{different hardware}.

Photonic neural networks perform matrix-vector multiplication using light
rather than electrical charge.
The key properties are:
(i) \emph{speed} --- computation completes at the speed of light traversing a
few millimetres of silicon;
(ii) \emph{energy} --- photons propagate through passive waveguides at
essentially zero energy cost;
(iii) \emph{parallelism} --- wavelength-division multiplexing (WDM) provides
free $N\times$ parallelism across attention heads.

The static compilation insight of this work is the following.
A transformer performing inference from a fixed, trained model never changes
its weights.
Therefore, the weight matrices can be \emph{permanently encoded} into the
physical geometry of a photonic chip --- the phase angles of its MZI
interferometers.
Inference then becomes pure physics: optical fields propagate through the chip,
and the computation is performed by light interference.
This is analogous to an application-specific integrated circuit (ASIC) versus a
GPU, but implemented in photons rather than electrons.

\paragraph{Contributions.} This paper presents:
\begin{enumerate}
  \item A complete open-source compiler (\S\ref{sec:compiler}) from MedGemma
        weights to MZI phase maps, validated at machine precision.
  \item A simulation framework (\S\ref{sec:simulation}) that applies the
        compiled Clements mesh to real MedGemma activations and compares the
        output to the digital ground truth.
  \item A rank sweep (\S\ref{sec:results}) from $r=64$ to full rank across all
        four layer-0 attention projections, with honest reporting of SVD
        approximation error, chip compilation fidelity, and hardware cost.
  \item A chip architecture specification (\S\ref{sec:architecture}) for
        220\,nm SOI silicon photonics with full power budget analysis.
\end{enumerate}

% ----------------------------------------------------------------
\section{Background}
\label{sec:background}

\subsection{MedGemma-4B-IT Architecture}

MedGemma-4B-IT is a multimodal vision-language model based on the Gemma-3
architecture~\cite{gemma3_2024}.
It comprises a SigLIP-So400M vision encoder (400\,M parameters) connected via
a linear projector to a Gemma-3-4B language model (3.6\,B parameters).
The language model has the following dimensions relevant to photonic compilation:

\begin{table}[h]
\centering
\caption{MedGemma-4B language model parameters relevant to photonic compilation.}
\label{tab:medgemma_arch}
\begin{tabular}{ll}
\toprule
\textbf{Parameter} & \textbf{Value} \\
\midrule
Hidden dimension $d_{\text{model}}$ & 2048 \\
Number of transformer layers & 46 \\
Number of attention heads & 8 \\
Number of key-value heads (GQA) & 4 \\
Per-head dimension & 256 \\
FFN inner dimension & 16{,}384 \\
\midrule
Q projection shape & $2048 \times 2560$ \\
K projection shape & $1024 \times 2560$ \\
V projection shape & $1024 \times 2560$ \\
O projection shape & $2560 \times 2048$ \\
\bottomrule
\end{tabular}
\end{table}

\subsection{Singular Value Decomposition for Photonic Compilation}

Any real or complex matrix $W \in \mathbb{R}^{m \times n}$ admits the singular
value decomposition (SVD):
\begin{equation}
  W = U \Sigma V^\dagger
  \label{eq:svd}
\end{equation}
where $U \in \mathbb{C}^{m \times m}$ and $V \in \mathbb{C}^{n \times n}$ are
unitary matrices and $\Sigma \in \mathbb{R}^{m \times n}$ is diagonal with
non-negative entries $\sigma_1 \geq \sigma_2 \geq \cdots \geq \sigma_k \geq 0$
($k = \min(m, n)$).

A rank-$r$ approximation is obtained by retaining only the $r$ largest singular
values:
\begin{equation}
  W_r = U_r \Sigma_r V_r^\dagger
  \label{eq:svd_truncated}
\end{equation}
where $U_r = U[:, :r]$, $\Sigma_r = \text{diag}(\sigma_1, \ldots, \sigma_r)$,
$V_r^\dagger = V^\dagger[:r, :]$.
The energy retained is $\sum_{i=1}^r \sigma_i^2 / \sum_{i=1}^k \sigma_i^2$.

Because $U_r$ and $V_r^\dagger$ are partial-unitary (their columns/rows are
orthonormal), each can be embedded into a full unitary and implemented as a
Clements MZI mesh.

\subsection{Clements Decomposition}
\label{sec:clements_background}

Clements et al.~\cite{clements2016} showed that any $N \times N$ unitary matrix
$U$ can be decomposed as a product of $2 \times 2$ MZI rotation matrices:
\begin{equation}
  U = D \cdot T_1 \cdot T_2 \cdots T_{N(N-1)/2}
  \label{eq:clements}
\end{equation}
where $D = \text{diag}(e^{i\psi_1}, \ldots, e^{i\psi_N})$ is a diagonal phase
screen and each $T_k$ acts on two adjacent modes $(m_k, m_k{+}1)$ with the
$2\times2$ transfer matrix:
\begin{equation}
  T(\theta, \varphi) =
  \begin{pmatrix}
    \cos(\theta/2) & -e^{-i\varphi}\sin(\theta/2) \\
    e^{i\varphi}\sin(\theta/2) & \cos(\theta/2)
  \end{pmatrix}.
  \label{eq:mzi_matrix}
\end{equation}
This decomposition uses exactly $N(N-1)/2$ MZIs --- provably optimal for
universal unitary synthesis.

For the forward pass, the unitary is applied as:
\begin{equation}
  U \mathbf{x} = (D \cdot T_1 \cdots T_M) \mathbf{x}
  = D \bigl( T_1 ( \cdots (T_M \mathbf{x}) \cdots ) \bigr),
  \label{eq:clements_forward}
\end{equation}
which is computed by first applying MZIs in \emph{descending column order}
(right-to-left elimination), then multiplying by the phase screen $D$.

% ----------------------------------------------------------------
\section{The PhotoMedGemma Compiler}
\label{sec:compiler}

\subsection{Pipeline Overview}

The compilation pipeline proceeds in six stages:

\begin{enumerate}
  \item \textbf{Model Parser}: loads MedGemma weights from HuggingFace
        (bfloat16, converted to float32 for numerical stability).
  \item \textbf{Layer Decomposer}: computes truncated SVD of each weight matrix
        at the selected rank $r$.
  \item \textbf{Clements Decomposer}: converts each unitary factor ($U_r$,
        $V_r^\dagger$) to $(r-1) \cdot r / 2$ MZI phase pairs $(\theta, \varphi)$
        plus a phase screen $\boldsymbol{\psi}$.
  \item \textbf{Phase Encoder}: quantizes phases to $B$-bit DAC resolution
        (default $B=12$ bits, giving $\lesssim 0.044^\circ$ phase error).
  \item \textbf{Netlist Generator}: outputs a photonic SPICE-like netlist
        (\texttt{.pntl} format).
  \item \textbf{GDS Generator}: produces a physical GDSII layout file for
        foundry submission using \texttt{gdsfactory}.
\end{enumerate}

\subsection{SVD Rank Selection}

Table~\ref{tab:rank_selection} summarises the rank selection strategy.
For the current proof-of-concept, we use $r = 64$.

\begin{table}[h]
\centering
\caption{Rank selection strategy for SVD-based photonic compilation.}
\label{tab:rank_selection}
\begin{tabular}{lccc}
\toprule
\textbf{Target} & \textbf{Energy threshold} & \textbf{Rank} & \textbf{MZI reduction} \\
\midrule
High accuracy      & 99.9\% & 256--512 & $16\times$ \\
Balanced           & 99\%   & 64--128  & $64\times$ \\
Aggressive         & 95\%   & 16--32   & $256\times$ \\
Demo / PoC         & 90\%   & 8--16    & $1024\times$ \\
\bottomrule
\end{tabular}
\end{table}

\subsection{Clements Convention and Numerical Precision}

The correct T-matrix convention for left-elimination Clements (as implemented
in \texttt{src/photonic/clements.py}) is:
\begin{equation}
  T(\theta, \varphi) =
  \begin{pmatrix}
    \cos(\theta/2) & -e^{-i\varphi}\sin(\theta/2) \\
    e^{i\varphi}\sin(\theta/2) & \cos(\theta/2)
  \end{pmatrix},
  \label{eq:T_correct}
\end{equation}
applied in \emph{descending column order}.
The compiled MZI phases in \texttt{output/real/phase\_map.json} reproduce both
the $V_r^\dagger$ and $U_r$ unitaries to errors $\varepsilon_U, \varepsilon_{V^\dagger}
\lesssim 8 \times 10^{-15}$ --- genuine machine precision --- for all four
layer-0 attention projections.

% ----------------------------------------------------------------
\section{Simulation and Validation}
\label{sec:simulation}

\subsection{Experimental Setup}

We ran MedGemma-4B-IT on five breast cancer PNG images
(\texttt{10025\_893612858.png}, \texttt{10011\_1031443799.png},
\texttt{10042\_102733848.png}, \texttt{10042\_1648588715.png},
\texttt{10042\_495770405.png})
using a Kaggle T4 GPU notebook with 4-bit NF4 quantization.
For each image, we captured:
\begin{itemize}
  \item The input hidden state $\mathbf{x} \in \mathbb{R}^{2560}$ into each
        layer-0 attention projection.
  \item The output $\mathbf{y}_{\text{MG}} = W_{\text{MG}} \mathbf{x}$
        produced by MedGemma's digital computation.
  \item The weight matrices $W_q, W_k, W_v, W_o$ (dequantized from NF4 to
        float32).
\end{itemize}
These arrays were downloaded as \texttt{.npy} files and used as ground truth
for the local photonic simulation.

\subsection{Error Metrics}

We report three independent error metrics:

\paragraph{Chip fidelity.}
\begin{equation}
  \varepsilon_{\text{chip}} =
  \frac{\| \text{MZI\_sim}(\mathbf{x}) - U_r V_r^\dagger \mathbf{x} \|}
       {\| U_r V_r^\dagger \mathbf{x} \|}
  \label{eq:chip_fidelity}
\end{equation}
This measures the floating-point precision of the Clements compilation.
Expected value: $\approx 10^{-15}$.

\paragraph{SVD approximation error.}
\begin{equation}
  \varepsilon_{\text{SVD}}(r) =
  \frac{\| W_r \mathbf{x} - W \mathbf{x} \|}
       {\| W \mathbf{x} \|}
  \label{eq:svd_error}
\end{equation}
This is the genuine information loss incurred by the rank-$r$ approximation.
It depends on the singular value spectrum of $W$ and the actual MedGemma input
$\mathbf{x}$.

\paragraph{Quantization error.}
\begin{equation}
  \varepsilon_{\text{quant}} =
  \frac{\| W_{\text{kaggle}} \mathbf{x} - \mathbf{y}_{\text{MG}} \|}
       {\| \mathbf{y}_{\text{MG}} \|}
  \label{eq:quant_error}
\end{equation}
This measures the error introduced by 4-bit NF4 weight quantization during
extraction from Kaggle.
Expected value: $< 0.3\%$.

\subsection{Validation Results}

Table~\ref{tab:comparison} reports all three metrics for rank-64 and full rank.

\begin{table}[h]
\centering
\caption{PhotoMedGemma validation on 5 real breast cancer images.
         Rank-64 chip vs.\ full MedGemma layer-0 attention.}
\label{tab:comparison}
\begin{tabular}{lcccccc}
\toprule
\multirow{2}{*}{\textbf{Projection}} &
\multirow{2}{*}{\textbf{Shape}} &
\multirow{2}{*}{\textbf{Chip fidelity}} &
\multicolumn{2}{c}{\textbf{SVD error}} &
\multirow{2}{*}{\textbf{Quant. err.}} \\
\cmidrule(lr){4-5}
& & & $r=64$ & $r=r_{\max}$ & \\
\midrule
\texttt{q\_proj} & $2048 \times 2560$ & $1.1 \times 10^{-15}$ & 0.186 & $1.9 \times 10^{-15}$ & 0.0017 \\
\texttt{k\_proj} & $1024 \times 2560$ & $1.2 \times 10^{-15}$ & 0.249 & $1.9 \times 10^{-15}$ & 0.0019 \\
\texttt{v\_proj} & $1024 \times 2560$ & $8.1 \times 10^{-16}$ & 0.669 & $3.1 \times 10^{-15}$ & 0.0022 \\
\texttt{o\_proj} & $2560 \times 2048$ & $1.1 \times 10^{-15}$ & 0.311 & $2.4 \times 10^{-15}$ & 0.0020 \\
\bottomrule
\end{tabular}
\end{table}

Three conclusions follow immediately from Table~\ref{tab:comparison}:
\begin{enumerate}
  \item \textbf{Compilation is exact.}
        Chip fidelity is at machine precision ($\approx 10^{-15}$) for all
        projections.
        The Clements algorithm faithfully encodes the target unitary.
  \item \textbf{SVD approximation error is rank-controlled.}
        At full rank, the photonic chip recovers $W\mathbf{x}$ to machine
        precision.
        At rank 64, the approximation error ranges from 0.19 to 0.67,
        reflecting the singular value spectra of individual weight matrices.
  \item \textbf{Quantization error is negligible.}
        The 4-bit NF4 extraction from Kaggle introduces $< 0.22\%$ error ---
        confirming that the weights used for compilation are accurate.
\end{enumerate}

% ----------------------------------------------------------------
\section{Rank Sweep: From 64 to Full Rank}
\label{sec:results}

We evaluate all four layer-0 attention projections at ranks
$r \in \{64, 128, 256, 512, 1024, 2048\}$ (where applicable).

\begin{table}[h]
\centering
\caption{Full rank sweep for \texttt{q\_proj} ($2048 \times 2560$, full rank = 2048).
         All chip-fidelity values are $\approx 10^{-15}$.}
\label{tab:rank_sweep_q}
\begin{tabular}{rrcrrr}
\toprule
\textbf{Rank} & \textbf{Energy} & \textbf{SVD error} & \textbf{MZIs/mesh} &
\textbf{Area (std)} & \textbf{Area (cmp)} \\
\midrule
    64 & 28.8\% & 0.186 &        2,016 &    4.9\,mm$^2$ &    1.6\,mm$^2$ \\
   128 & 38.2\% & 0.166 &        8,128 &   19.7\,mm$^2$ &    6.6\,mm$^2$ \\
   256 & 52.8\% & 0.142 &       32,640 &   78.6\,mm$^2$ &   26.2\,mm$^2$ \\
   512 & 72.6\% & 0.108 &      130,816 &  314.6\,mm$^2$ &  104.9\,mm$^2$ \\
  1024 & 92.3\% & 0.061 &      523,776 & 1258.3\,mm$^2$ &  419.4\,mm$^2$ \\
  2048 & 100\%  & 0.000 &    2,096,128 & 5033.2\,mm$^2$ & 1677.7\,mm$^2$ \\
\bottomrule
\end{tabular}
\end{table}

\noindent
\textit{Standard Si}: 150\,\textmu m stage $\times$ 8\,\textmu m pitch per MZI.
\textit{Compact Si}: 80\,\textmu m stage $\times$ 5\,\textmu m pitch per MZI.
Each projection requires 2 unitary meshes (Vh + U) plus a sigma (VOA) stage.

\begin{table}[h]
\centering
\caption{SVD error vs.\ rank for all four layer-0 attention projections.}
\label{tab:rank_sweep_all}
\begin{tabular}{rcccc}
\toprule
\textbf{Rank} & \textbf{q\_proj} & \textbf{k\_proj} & \textbf{v\_proj} & \textbf{o\_proj} \\
\midrule
    64 & 0.186 & 0.249 & 0.669 & 0.311 \\
   128 & 0.166 & 0.223 & 0.555 & 0.282 \\
   256 & 0.142 & 0.178 & 0.405 & 0.241 \\
   512 & 0.108 & 0.115 & 0.229 & 0.169 \\
  1024 & 0.061 & 0.000 & 0.000 & 0.084 \\
  2048 & 0.000 & ---   & ---   & 0.000 \\
\bottomrule
\end{tabular}
\end{table}

The v\_proj matrix has the slowest singular value decay, indicating that
it requires higher rank for a given approximation quality.
This is a physical property of the weight matrix, not a limitation of the
compiler.

% ----------------------------------------------------------------
\section{Chip Architecture}
\label{sec:architecture}

\subsection{Platform: 220\,nm Silicon-on-Insulator (SOI)}

We target the 220\,nm SOI process available at imec (iSiPP50G), AIM Photonics,
and GlobalFoundries (SiPh 45CLO).
Key parameters:
\begin{itemize}
  \item Waveguide: 450\,nm wide $\times$ 220\,nm tall Si core, SiO$_2$ cladding
  \item Wavelength: 1310\,nm (O-band, low dispersion)
  \item Phase shifter: thermal TiN heater, $\Delta\phi = \pi$ at 10--20\,mW
  \item Propagation loss: 2--3\,dB/cm
  \item Coupling loss: 2--4\,dB per grating coupler
  \item Detector: Ge photodiode, $> 40$\,GHz bandwidth, 0.9\,A/W responsivity
\end{itemize}

\subsection{Single Chip: 10\,mm $\times$ 10\,mm}

Each chip implements one $64 \times 64$ Clements unitary mesh (2,016 MZIs)
in an 8\,mm $\times$ 8\,mm active area, with 64 grating couplers on each edge
for fiber I/O.
The heater control routing uses the top metal layer: 1,024 heater lines routed
to bond pads connected to DAC chips on the PCB.

\subsection{Power Budget}

\begin{table}[h]
\centering
\caption{Power budget per transformer layer (rank-64, 46 layers).}
\label{tab:power}
\begin{tabular}{lrl}
\toprule
\textbf{Component} & \textbf{Power} & \textbf{Notes} \\
\midrule
Laser sources (8$\times$ WDM) &  800\,mW & $\approx 100$\,mW per wavelength \\
Phase shifters (static)        &    0\,mW & Thermal; no steady-state draw \\
Photodetectors                 &   50\,mW & Ge PD array \\
Electronic control (FPGA)      &    5\,W  & Softmax, LayerNorm, routing \\
Memory (KV cache)              &    2\,W  & LPDDR5 \\
\midrule
\textbf{Total}                 & $\approx 8$\,W & \textbf{vs.\ $\approx 300$\,W for A100 GPU} \\
\bottomrule
\end{tabular}
\end{table}

At 1 token per millisecond: $8\,\text{mJ/token}$ versus
$300\,\text{mJ/token}$ for an A100 --- a \textbf{37$\times$ improvement in
energy efficiency}.

% ----------------------------------------------------------------
\section{Discussion}
\label{sec:discussion}

\paragraph{On SVD approximation error.}
The rank-64 SVD errors (0.19--0.67) represent genuine information loss relative
to the full weight matrix.
They do \emph{not} indicate a flaw in the photonic compiler.
The chip executes the rank-64 approximation to machine precision.
The question of how much task accuracy is preserved at rank 64 versus full rank
requires end-to-end evaluation on medical benchmarks (MedQA, PubMedQA), which
is the subject of ongoing work.

\paragraph{On chip count.}
A rank-64 implementation of all four layer-0 projections requires
$4 \times 2$ unitary meshes = 8 chips per layer.
Scaling to all 46 layers and 7 weight matrices per layer: $\approx 2{,}576$
chips at rank-64.
This is large but manufacturable with a multi-chip module approach at current
silicon photonics scales.

\paragraph{On fabrication.}
The compilation output (\texttt{output/real/phase\_map.json},
\texttt{output/real/netlist/}, \texttt{output/real/gds/}) is ready for
submission to a photonic foundry.
The GDS layout uses standard \texttt{gdsfactory} cells compatible with 220\,nm
SOI PDKs from imec, AIM Photonics, and GlobalFoundries.

% ----------------------------------------------------------------
\section{Conclusion}
\label{sec:conclusion}

We have presented PhotoMedGemma: a complete compiler and simulation framework
that translates MedGemma-4B-IT's attention weights into Clements MZI phase
maps for silicon photonic chips.
The compilation achieves machine precision ($\approx 7 \times 10^{-15}$) for
both the $V_r^\dagger$ and $U_r$ unitaries.
The rank-$r$ approximation quality is a continuous design parameter: at full
rank the chip recovers the digital computation exactly; at rank 64 it trades
hardware simplicity for approximation error.

The project is motivated by a concrete human need: deploying clinical-grade
medical AI at Kenyatta National Hospital and hospitals throughout the global
south, without cloud dependency, data sovereignty risk, or energy crisis.
The photonic hardware projects $37\times$ better energy efficiency than current
GPU infrastructure.

All code, compiled phase maps, and experimental results are open-source.

% ----------------------------------------------------------------
\bibliographystyle{plain}
\begin{thebibliography}{9}

\bibitem{medgemma2024}
Google DeepMind,
\textit{MedGemma: A Family of Medical AI Models},
Technical Report, 2024.

\bibitem{gemma3_2024}
Google DeepMind,
\textit{Gemma 3 Technical Report},
arXiv:2408.00118, 2024.

\bibitem{clements2016}
W.R.\ Clements, P.C.\ Humphreys, B.J.\ Metcalf, W.S.\ Kolthammer,
and I.A.\ Walmsley,
``Optimal design for universal multiport interferometers,''
\textit{Optica}, 3(12):1460--1465, 2016.

\bibitem{reck1994}
M.\ Reck, A.\ Zeilinger, H.J.\ Bernstein, and P.\ Bertani,
``Experimental realization of any discrete unitary operator,''
\textit{Physical Review Letters}, 73(1):58, 1994.

\bibitem{shen2017}
Y.\ Shen, N.C.\ Harris, S.\ Skirlo, M.\ Prabhu, T.\ Baehr-Jones,
M.\ Hochberg, X.\ Sun, S.\ Zhao, H.\ Larochelle, D.\ Englund, and M.\ Solja\v{c}i\'{c},
``Deep learning with coherent nanophotonic circuits,''
\textit{Nature Photonics}, 11(7):441--446, 2017.

\bibitem{bandyopadhyay2022}
S.\ Bandyopadhyay, A.\ Sludds, S.\ Krastanov, R.\ Hamerly,
N.\ Harris, D.\ Bunandar, M.\ Streshinsky, M.\ Hochberg, and D.\ Englund,
``Single chip photonic deep neural network with accelerated training,''
\textit{arXiv:2208.01623}, 2022.

\end{thebibliography}

\end{document}
